\subsection{PRODUCT OF ALTERNATING MULTILINEAR FORMS}

Let A be a commutative ring and M an A-module. For each integer \(r \geq 0\) , denote by \(\operatorname{Alt}^{r}(\mathbf{M})\) the A-module of alternating r-linear forms on M; it can be identified with the dual of the A-module \(\Lambda^{r}(\mathbf{M})\) (Algebra, Chap. III, §7, no. 4, Prop. 7). Let \(u \in \operatorname{Alt}^{s}(\mathbf{M})\) and \(v \in \operatorname{Alt}^{r}(\mathbf{M})\) ; recall (Algebra, Chap. III, §11, no. 2, Example 3) that the alternating product of u and v is the element \(u \wedge v \in \operatorname{Alt}^{s+r}(\mathbf{M})\) defined by

\[
(u \wedge v) (x _ {1}, \dots , x _ {s + r}) = \sum_ {\sigma \in \mathfrak {S} _ {s, r}} \varepsilon_ {\sigma} u (x _ {\sigma (1)}, \dots , x _ {\sigma (s)}) v (x _ {\sigma (s + 1)}, \dots , x _ {\sigma (s + r)}),
\]

where \(S_{s,r}\) is the subset of the symmetric group \(S_{s+r}\) consisting of the permutations whose restrictions to \((1,s)\) and \((s+1,s+r)\) are increasing.

Now let

\[
0 \longrightarrow \mathrm{M} ^ {\prime} \stackrel {i} {\longrightarrow} \mathrm{M} \stackrel {p} {\longrightarrow} \mathrm{M} ^ {\prime \prime} \longrightarrow 0
\]

be an exact sequence of free A-modules, of ranks \(r, r + s\) and \(s\) , respectively.

Lemma 1. There exists an A-bilinear map from \(\mathrm{Alt}^s (\mathbf{M}'')\times \mathrm{Alt}^r (\mathbf{M}')\) to \(\mathrm{Alt}^{s + r}(\mathbf{M})\) , denoted by \((u,v)\mapsto u\cap v\) , and characterized by either of the following two properties:

\begin{enumerate}
\def\labelenumi{\alph{enumi})}
\item
  Denote by \(u_{1} \in \operatorname{Alt}^{s}(\operatorname{M})\) the form \((x_{1}, \ldots, x_{s}) \mapsto u(p(x_{1}), \ldots, p(x_{s}))\) , and let \(v_{1} \in \operatorname{Alt}^{r}(\operatorname{M})\) be a form such that \(v_{1}(i(x_{1}^{\prime}), \ldots, i(x_{r}^{\prime})) = v(x_{1}^{\prime}, \ldots, x_{r}^{\prime})\) for \(x_{1}^{\prime}, \ldots, x_{r}^{\prime}\) in \(M^{\prime}\) ; then \(u \cap v = u_{1} \wedge v_{1}\) .
\item
  For all \(x_{1},\ldots,x_{s}\) in M and \(x_{1}^{\prime},\ldots,x_{r}^{\prime}\) in \(M^{\prime}\) ,
\end{enumerate}

\[
(u \cap v) (x _ {1}, \dots , x _ {s}, i (x _ {1} ^ {\prime}), \dots , i (x _ {r} ^ {\prime})) = u (p (x _ {1}), \dots , p (x _ {s})) v (x _ {1} ^ {\prime}, \dots , x _ {r} ^ {\prime}).\tag{1}
\]

The map \(\varphi : \mathrm{Alt}^s (\mathrm{M}'')\otimes_{\mathrm{A}}\mathrm{Alt}^r (\mathrm{M}')\to \mathrm{Alt}^{s + r}(\mathrm{M})\) such that \(\varphi (u\otimes v) = u\cap v\) is an isomorphism of free A-modules of rank one.

The existence of a form \(v_{1}\) satisfying condition \(a\) ) follows from the fact that \(\Lambda^r (i)\) induces an isomorphism from \(\Lambda^r (\mathrm{M}')\) to a direct factor submodule of \(\Lambda^r (\mathrm{M})\) (Algebra, Chap. III, §7, no. 2). Let \(v_{1}\) be such a form; put \(u\cap v = u_{1}\wedge v_{1}\) . Formula (1) is then satisfied, since if we put \(i(x_k') = x_{s + k}\) for \(1\leq k\leq r\) , the only element \(\sigma\) of \(\mathfrak{S}_{s,r}\) such that \(p(x_{\sigma (i)})\neq 0\) for \(1\leq i\leq s\) is the identity permutation. On the other hand, formula (1) determines \(u\cap v\) uniquely: indeed, let \((e_1',\ldots ,e_r')\) be a basis of \(\mathrm{M}'\) , \((f_1'',\dots,f_s'')\) a basis of \(\mathrm{M}''\) , and \(f_{1},\ldots ,f_{s}\) elements of \(\mathrm{M}\) such that \(p(f_i) = f_i''\) for \(1\leq i\leq s\) . Then \((f_{1},\ldots ,f_{s},i(e_1'),\ldots ,i(e_r'))\) is a basis of \(\mathrm{M}\) (Algebra, Chap. II, §1, no. 11, Prop. 21), and formula (1) can be written

\[
(u \cap v) (f _ {1}, \dots , f _ {s}, i (e _ {1} ^ {\prime}), \dots , i (e _ {r} ^ {\prime})) = u (f _ {1} ^ {\prime \prime}, \dots , f _ {s} ^ {\prime \prime}) v (e _ {1} ^ {\prime}, \dots , e _ {r} ^ {\prime});\tag{2}
\]

but an element of \(\mathrm{Alt}^{s + r}(\mathbf{M})\) is determined by its value on a basis.

It follows from the preceding that each of the conditions a) and b) determines the product \(u \cap v\) uniquely; it is clear that this product is bilinear. Finally, the last assertion of the lemma follows from formula (2).
% semantic-fragments: legacy-input-seam
\relax{}
